\documentclass{article}
\usepackage{todonotes}
\usepackage{tcolorbox}
\usepackage{xcolor} 
\usepackage{amsmath}
\usepackage{amssymb}

\newtheorem{test}{Test}[section]
\newtheorem{bsp}[test]{Beispiel}
\newtheorem{info}[test]{Zusatzinformation}
\newtheorem{reg}[test]{Regel}
\newtheorem{satz}[test]{Satz}
\newtheorem{lem}[test]{Lemma}

% Definition der tcolorbox für 'lem' (Lemma)
\tcolorboxenvironment{lem}{
  colback=white,
  colframe=black,
  boxrule=0.8pt,
  arc=2mm,
  top=1mm,
  bottom=1mm,
  left=1mm,
  right=1mm,
  parbox=false, % Keine Parbox (keine Verformung der Box)
  enhanced
}

\newtheorem{defi}[test]{Definition}
% Definition der tcolorbox für 'defi' (Definition)
\tcolorboxenvironment{defi}{
  colback=white,
  colframe=black,
  boxrule=0.8pt,
  arc=2mm,
  top=1mm,
  bottom=1mm,
  left=1mm,
  right=1mm,
  enhanced
}

% Definition der Randnotizen
\newcommand{\bluenote}[1]{\todo[bordercolor=blue, backgroundcolor=white, textcolor=blue, size=\small]{\textcolor{blue}{#1}}}
\newcommand{\orangenote}[1]{\todo[bordercolor=orange, backgroundcolor=white, textcolor=orange, size=\small]{\textcolor{orange}{#1}}}

\begin{document}
  
\section{Beispiel mit tcolorbox und Randnotizen außerhalb der Boxen}

% Eine Beispieldefinition
\begin{defi}
  Eine \textbf{Definition} ist eine präzise und klare Aussage, die den Begriff beschreibt.
  % Randnotiz innerhalb der Box, aber außerhalb der Box angezeigt
  \bluenote{Denke daran, nach der Definition ein Beispiel zu geben!}
\end{defi}

% Eine Beispielregel mit Randnotizen
\begin{defi}
  Eine \textbf{Regel} ist eine vorgeschriebene Handlung oder Bedingung.
  \orangenote{Überprüfe, ob du alle Bedingungen berücksichtigst.}
\end{defi}

\end{document}
